\chapter{Cơ sở lý thuyết và các công trình nghiên cứu liên quan}
\label{Chapter2}

Chương này tổng hợp các hướng nghiên cứu có liên quan trực tiếp đến bài toán phân loại ảnh X-quang xương đa phương thức được xem xét trong luận văn. Thay vì khảo sát toàn bộ các mô hình y sinh đa phương thức, nội dung được tổ chức theo các quyết định thiết kế của hệ thống: lựa chọn nguồn văn bản, kế thừa mô hình ngôn ngữ -- thị giác y khoa, bảo toàn thông tin trên ảnh độ phân giải cao, tinh chỉnh hiệu quả tham số (Parameter-Efficient Fine-Tuning -- PEFT), dung hợp hai phương thức, phân loại trong điều kiện mất cân bằng, phát hiện dữ liệu ngoài phân phối (Out-of-Distribution -- OOD) và phân tích bản đồ chú ý. Cách tổ chức này nhằm làm rõ mối liên hệ giữa các công trình trước đây, khoảng trống còn lại và quy trình (pipeline) được trình bày ở Chương 3.

Các kết quả định lượng trong những nghiên cứu được trích dẫn không được đặt cạnh nhau như một bảng xếp hạng, do khác biệt về bộ dữ liệu, cách chia tập, loại bệnh lý và chỉ số đánh giá. Thay vào đó, các bảng trong chương tập trung so sánh giả định, loại đầu vào, cơ chế học và mức độ liên quan với bài toán của luận văn.

\section{Học đa phương thức trong chẩn đoán hình ảnh y khoa}

\subsection{Tính bổ sung giữa hình ảnh và thông tin lâm sàng}

Ảnh X-quang mô tả cấu trúc giải phẫu và các biểu hiện có thể quan sát được tại thời điểm chụp, trong khi bệnh sử lâm sàng cung cấp thêm bối cảnh như triệu chứng, thời gian khởi phát, tuổi và vị trí đau. Hai nguồn dữ liệu vì vậy có thể mang thông tin bổ sung cho nhau. Tuy nhiên, mức độ hữu ích của văn bản phụ thuộc đáng kể vào thời điểm và mục đích thu nhận. Bệnh sử được ghi nhận trước khi đọc ảnh chủ yếu phản ánh bối cảnh lâm sàng; ngược lại, báo cáo X-quang có thể chứa mô tả tổn thương hoặc kết luận chẩn đoán được dùng làm nhãn. Nếu báo cáo này xuất hiện ở đầu vào, mô hình có thể dựa vào từ khóa chẩn đoán thay vì học quan hệ giữa ảnh, bệnh sử và bệnh lý. Đây là một dạng rò rỉ thông tin đích cần được kiểm soát ở cả huấn luyện và suy luận.

Các nghiên cứu trước cho thấy nhiều mức độ kết hợp ảnh và dữ liệu lâm sàng. Mei và cộng sự \cite{XMei} kết hợp đặc trưng CT ngực với thông tin lâm sàng bằng một mô-đun dự đoán chung. TieNet \cite{TieNet} sử dụng cơ chế chú ý để liên kết từ trong báo cáo với vùng ảnh X-quang ngực. IRENE \cite{IRENE} biểu diễn nhiều loại dữ liệu lâm sàng dưới dạng token và cho phép chúng tương tác trong Transformer. FTB \cite{FTB} giữ cố định các backbone ảnh và văn bản, đồng thời học mô-đun kết nối giữa hai phía. Các công trình này cung cấp bằng chứng phương pháp luận cho việc khai thác tương tác đa phương thức, nhưng loại văn bản, miền giải phẫu và mục tiêu dự đoán của chúng không hoàn toàn trùng với bài toán u xương trong luận văn. Bảng \ref{tab:ch2_multimodal_clinical} tóm tắt sự khác biệt chính giữa các hướng tiếp cận này.

\begin{table}[htbp]
\centering
\small
\renewcommand{\arraystretch}{1.2}
\begin{tabularx}{\textwidth}{L{2.1cm} L{2.6cm} L{2.8cm} X X}
\toprule
\textbf{Công trình} & \textbf{Dữ liệu hình ảnh} & \textbf{Nguồn văn bản / dữ liệu phụ} & \textbf{Cơ chế kết hợp chính} & \textbf{Liên hệ và giới hạn đối với luận văn} \\
\midrule
Mei và cộng sự \cite{XMei} & CT ngực & Triệu chứng và thông tin lâm sàng & Ghép đặc trưng ảnh và lâm sàng cho dự đoán chung & Hỗ trợ giả định về tính bổ sung của dữ liệu lâm sàng; khác miền ảnh và bệnh lý. \\
\addlinespace
TieNet \cite{TieNet} & X-quang ngực & Báo cáo X-quang & Chú ý giữa từ và vùng ảnh & Cung cấp tiền lệ cho tương tác cục bộ ảnh--văn bản; nguồn văn bản có thể chứa thông tin gần với nhãn. \\
\addlinespace
IRENE \cite{IRENE} & Nhiều loại ảnh y khoa & Dữ liệu lâm sàng không đồng nhất & Mã hóa thống nhất bằng token và Transformer & Liên quan đến tương tác token đa phương thức; kiến trúc và phạm vi dữ liệu rộng hơn bài toán hiện tại. \\
\addlinespace
FTB \cite{FTB} & Ảnh y khoa & Văn bản y khoa & Giữ cố định backbone và học khối kết nối & Gợi ý cách thích nghi tiết kiệm tham số; không sử dụng cùng bộ phân loại và quy trình OOD của luận văn. \\
\bottomrule
\end{tabularx}
\caption{Tổng hợp một số hướng kết hợp hình ảnh và dữ liệu lâm sàng.}
\label{tab:ch2_multimodal_clinical}
\end{table}

\subsection{Lựa chọn nguồn văn bản và kiểm soát rò rỉ nhãn}

Từ góc độ thiết kế thí nghiệm, cần phân biệt rõ \textit{bệnh sử lâm sàng được ghi nhận trước khi có kết luận đọc ảnh} với \textit{báo cáo đọc ảnh}. Hai nguồn này có thể cùng được lưu dưới dạng văn bản nhưng mang vai trò khác nhau trong quy trình khám. Với mục tiêu dự đoán loại bệnh từ ảnh và thông tin sẵn có trước khi có kết luận X-quang, bệnh sử lâm sàng phù hợp hơn với điều kiện suy luận dự kiến của đề tài. Việc sử dụng cùng một loại văn bản trong cả hai pha huấn luyện còn giúp hạn chế sai khác phân phối đầu vào giữa pha căn chỉnh và pha phân loại.

Do đó, luận văn sử dụng ảnh X-quang và bệnh sử lâm sàng trong pipeline chính. Báo cáo đọc ảnh không được xem là một phương thức đầu vào tương đương, mà chỉ phù hợp cho các phân tích riêng về nguy cơ rò rỉ nhãn. Lựa chọn này không hàm ý rằng báo cáo X-quang không hữu ích trong các nhiệm vụ khác, chẳng hạn sinh báo cáo hoặc truy hồi ảnh--văn bản; nó chỉ phản ánh phạm vi dự đoán và điều kiện đánh giá của đề tài.

\section{Mô hình ngôn ngữ -- thị giác y khoa}

\subsection{Học biểu diễn tương phản ảnh--văn bản}

Các mô hình ngôn ngữ -- thị giác y khoa thường sử dụng một bộ mã hóa ảnh và một bộ mã hóa văn bản, sau đó đưa hai loại biểu diễn về cùng không gian để tối ưu một mục tiêu căn chỉnh. Với một mini-batch gồm các cặp ảnh và văn bản, mục tiêu tương phản làm tăng độ tương đồng của các cặp liên quan và giảm độ tương đồng của các cặp còn lại. Cơ chế này cho phép tận dụng văn bản đi kèm ảnh trong giai đoạn tiền huấn luyện hoặc thích nghi miền mà không nhất thiết phải xây dựng một bộ phân loại riêng cho từng nhãn.

ConVIRT \cite{ConVIRT} cho thấy khả năng học biểu diễn ảnh y khoa từ các cặp ảnh--văn bản bằng mục tiêu tương phản hai chiều. GLoRIA \cite{GLoRIA} mở rộng căn chỉnh toàn cục bằng quan hệ cục bộ giữa vùng ảnh và từ. MedCLIP \cite{MedCLIP} xem xét vấn đề cặp âm tính giả và đề xuất Semantic Matching để sử dụng quan hệ ngữ nghĩa mềm thay cho giả định mỗi ảnh chỉ khớp với đúng một văn bản. BiomedCLIP \cite{BiomedCLIP} được tiền huấn luyện trên PMC-15M với ViT và PubMedBERT, qua đó cung cấp một điểm khởi tạo đa phương thức thuộc miền y sinh.

Trong dữ liệu bệnh lý, nhiều bệnh nhân cùng lớp có thể có các mô tả khác nhau nhưng vẫn chia sẻ nội dung ngữ nghĩa liên quan. Nếu mọi cặp không trùng chỉ số trong mini-batch đều bị xem là âm tính, mục tiêu một--một có thể tạo ra tín hiệu tối ưu chưa phù hợp. Semantic Matching là tiền đề gần với pha 1 của luận văn, trong đó các cặp cùng lớp được gán quan hệ mềm sau khi chuẩn hóa theo hàng, còn cặp ảnh--bệnh sử tương ứng vẫn giữ mức khớp cao nhất. Bảng \ref{tab:ch2_medvlm} đối chiếu các mô hình theo bộ mã hóa, mục tiêu học và vai trò đối với pipeline.

\begin{table}[htbp]
\centering
\small
\renewcommand{\arraystretch}{1.2}
\begin{tabularx}{\textwidth}{L{2.2cm} L{2.5cm} L{2.5cm} X X}
\toprule
\textbf{Mô hình} & \textbf{Bộ mã hóa ảnh} & \textbf{Bộ mã hóa văn bản} & \textbf{Mục tiêu hoặc cơ chế chính} & \textbf{Vai trò đối với pipeline đề xuất} \\
\midrule
ConVIRT \cite{ConVIRT} & ResNet-50 & BERT & Tương phản ảnh--văn bản hai chiều & Cơ sở cho học biểu diễn y khoa từ dữ liệu ghép cặp. \\
\addlinespace
GLoRIA \cite{GLoRIA} & ResNet-50 & BioClinicalBERT & Căn chỉnh toàn cục và cục bộ & Gợi ý rằng thông tin vùng ảnh và token văn bản có thể bổ sung cho biểu diễn toàn cục. \\
\addlinespace
MedCLIP \cite{MedCLIP} & ResNet-50 hoặc Swin Transformer & BioClinicalBERT & Semantic Matching và dữ liệu không ghép cặp hoàn toàn & Cơ sở trực tiếp cho mục tiêu soft-target ở pha 1. \\
\addlinespace
BiomedCLIP \cite{BiomedCLIP} & ViT-Base & PubMedBERT & Tiền huấn luyện tương phản trên PMC-15M & Backbone y sinh được kế thừa và thích nghi cho ảnh X-quang xương. \\
\bottomrule
\end{tabularx}
\caption{So sánh các mô hình ngôn ngữ -- thị giác y khoa có liên quan trực tiếp.}
\label{tab:ch2_medvlm}
\end{table}

\subsection{Lý do kế thừa mô hình nền tảng thay vì huấn luyện từ đầu}

Huấn luyện một mô hình đa phương thức từ đầu thường cần số lượng lớn cặp ảnh--văn bản và tài nguyên tính toán đáng kể. Trong bài toán u xương, số mẫu ở từng lớp có thể hạn chế và phân bố lớp không đồng đều. Kế thừa BiomedCLIP giúp bắt đầu từ các bộ mã hóa đã có khả năng biểu diễn ngôn ngữ và hình ảnh y sinh, sau đó chỉ thích nghi các thành phần cần thiết cho miền X-quang xương. Tuy vậy, tiền huấn luyện trên dữ liệu y sinh rộng không bảo đảm mô hình đã phù hợp với phân bố đích; vì vậy, thích nghi miền và đánh giá bằng các đối chứng vẫn là cần thiết.

Các mô hình đa phương thức sinh văn bản quy mô lớn phù hợp hơn với hội thoại, hỏi--đáp hoặc sinh báo cáo. Do luận văn tập trung vào phân loại đóng, phát hiện OOD và phân tích attention, nhóm mô hình này không được khảo sát như một nhánh chính. Flamingo \cite{Alayrac2022Flamingo} chỉ được nhắc lại ở phần sau như một ví dụ về cơ chế gom token thị giác, thay vì làm nền tảng cho tác vụ sinh văn bản.

\section{Biểu diễn ảnh độ phân giải cao và thông tin không gian}

\subsection{Giới hạn của phép co giãn ảnh toàn cục}

Các backbone ViT được tiền huấn luyện thường sử dụng kích thước đầu vào cố định. Khi toàn bộ ảnh X-quang có kích thước lớn được co về một ảnh nhỏ, cấu trúc tổng thể vẫn được duy trì nhưng một số chi tiết cục bộ có thể bị làm mờ. Ngược lại, đưa trực tiếp ảnh độ phân giải cao vào Transformer làm tăng số token và chi phí attention. FPT \cite{FPT} xem xét bài toán chuyển giao trên ảnh y khoa độ phân giải cao và sử dụng lựa chọn token cùng mạng nhánh để giảm yêu cầu bộ nhớ. Mặc dù cơ chế của FPT khác pipeline trong luận văn, công trình này cho thấy cần xem xét đồng thời độ chi tiết và chi phí tính toán khi thích nghi mô hình nền tảng cho ảnh y khoa.

Một hướng khác là chuyển tập đặc trưng đầu vào lớn thành ngân sách latent cố định. Perceiver \cite{Jaegle2021Perceiver} sử dụng attention bất đối xứng để ánh xạ đầu vào có số phần tử lớn vào một tập latent nhỏ. Flamingo \cite{Alayrac2022Flamingo} sử dụng Perceiver Resampler để tạo số token thị giác cố định trước khi kết nối với mô hình ngôn ngữ. Các nghiên cứu này cung cấp cơ sở khái niệm cho việc kiểm soát ngân sách token, nhưng không đồng nghĩa với việc pipeline của luận văn sử dụng learned queries giống Perceiver.

\subsection{Kết hợp ngữ cảnh toàn cục và vùng cục bộ}

Đối với ảnh X-quang xương, nhánh toàn cục hỗ trợ duy trì vị trí giải phẫu và cấu trúc chung, trong khi các vùng cục bộ giữ lại thông tin chi tiết ở độ phân giải đầu vào của backbone. Tuy nhiên, nếu chỉ gom trung bình các token cục bộ, mô hình có thể mất thông tin về vị trí tương đối của từng vùng. Vì vậy, tọa độ đã chuẩn hóa có thể được bổ sung vào token trước khi dung hợp, với hệ số điều khiển để thông tin vị trí không lấn át nội dung ảnh.

Pipeline cuối cùng sử dụng một ảnh toàn cục (\textit{global view}) và bốn ảnh cục bộ (\textit{local view}) được chọn theo chiến lược chọn thưa có trọng tâm (\textit{sparse-focal}) chỉ từ tín hiệu ảnh. Mỗi vùng cục bộ đóng góp CLS token và token tóm tắt patch; tám token được truyền trực tiếp qua mô-đun không gian có bổ sung tọa độ. Cách thiết kế này gần với một ngân sách token xác định hơn là một bộ resampler dùng truy vấn học được. Hiệu quả của toàn bộ nhánh ảnh độ phân giải cao cần được kiểm chứng bằng thí nghiệm loại bỏ thành phần, thay vì được giả định từ các công trình dùng kiến trúc khác. Nếu cần tách riêng đóng góp của tọa độ, quy trình thực nghiệm phải bổ sung một cấu hình chỉ tắt thành phần này trong khi giữ nguyên cách chọn vùng. Những đánh đổi chính được tổng hợp trong Bảng \ref{tab:ch2_high_resolution}.

\begin{table}[htbp]
\centering
\small
\renewcommand{\arraystretch}{1.2}
\begin{tabularx}{\textwidth}{L{2.6cm} L{3.0cm} X X}
\toprule
\textbf{Chiến lược} & \textbf{Đặc điểm} & \textbf{Khả năng đáp ứng} & \textbf{Hạn chế cần xem xét} \\
\midrule
Co giãn ảnh toàn cục & Một ảnh được đưa về kích thước cố định & Chi phí thấp, duy trì bố cục tổng thể & Có thể làm giảm chi tiết cục bộ nhỏ. \\
\addlinespace
Xử lý trực tiếp ảnh độ phân giải cao & Tăng kích thước ảnh hoặc số patch & Giữ nhiều chi tiết hơn & Tăng bộ nhớ và chi phí attention; FPT \cite{FPT} xem xét sự đánh đổi này. \\
\addlinespace
Latent resampling & Nén tập token đầu vào vào số latent cố định & Kiểm soát độ dài chuỗi & Hiệu quả phụ thuộc dữ liệu và cách tối ưu truy vấn học được; đại diện bởi Perceiver \cite{Jaegle2021Perceiver} và Flamingo \cite{Alayrac2022Flamingo}. \\
\addlinespace
Global--local với token cố định & Kết hợp ảnh toàn cục, các vùng cục bộ và tọa độ & Cân bằng ngữ cảnh, chi tiết và ngân sách token & Phụ thuộc vào quy tắc chọn vùng; cần đánh giá bằng ablation trên cùng dữ liệu. \\
\bottomrule
\end{tabularx}
\caption{Các chiến lược xử lý ảnh độ phân giải cao và sự đánh đổi chính.}
\label{tab:ch2_high_resolution}
\end{table}

\section{Tinh chỉnh hiệu quả tham số}

\subsection{Các nhóm phương pháp}

Tinh chỉnh toàn bộ tham số cho phép mô hình thay đổi linh hoạt theo dữ liệu đích, nhưng yêu cầu lưu gradient và trạng thái optimizer cho toàn bộ mô hình nền (backbone). Dò tuyến tính (\textit{linear probing}) giảm chi phí bằng cách chỉ học đầu phân loại, song giới hạn khả năng thích nghi biểu diễn. Các phương pháp PEFT nằm giữa hai lựa chọn này: phần lớn trọng số tiền huấn luyện được giữ cố định, còn một tập tham số nhỏ được thêm vào hoặc lựa chọn để cập nhật.

Prompt tuning học các token đầu vào hoặc prompt trung gian; PLIP \cite{PLIP}, BiomedCoOp \cite{BiomedCoOp} và FPT \cite{FPT} là các ví dụ trong ảnh y khoa. Adapter tuning chèn các mô-đun nhỏ vào mạng hoặc giữa các backbone; FTB \cite{FTB} đại diện cho hướng học khối kết nối đa phương thức. Những nghiên cứu này có giá trị so sánh, nhưng pipeline của luận văn không triển khai prompt tuning làm thành phần chính.

\subsection{Thích nghi hạng thấp}

LoRA \cite{Hu2022LoRA} giữ cố định ma trận trọng số tiền huấn luyện $W$ và biểu diễn phần cập nhật bằng tích hai ma trận hạng thấp:
\begin{equation}
    W' = W + \Delta W, \qquad \Delta W = BA,
\end{equation}
trong đó $A \in \mathbb{R}^{r\times k}$, $B \in \mathbb{R}^{d\times r}$ và $r \ll \min(d,k)$. Cách tham số hóa này giảm số tham số được tối ưu so với cập nhật toàn bộ $W$. LCBB \cite{LCBB} khảo sát PEFT trên các mô hình nền tảng thị giác y khoa và cung cấp thêm bằng chứng thực nghiệm cho việc sử dụng LoRA trong miền y tế.

LoRA khác với selective tuning thuần túy vì nó bổ sung các ma trận có thể học, thay vì chỉ chọn một tập con trọng số gốc. LoRA cũng khác QLoRA ở chỗ QLoRA kết hợp thích nghi hạng thấp với lượng tử hóa trọng số nền. Pipeline của luận văn sử dụng LoRA không lượng tử hóa và không khảo sát QLoRA, nên không đưa ra kết luận thực nghiệm về ảnh hưởng của lượng tử hóa. Các adapter được giữ tách rời sau pha 1 để có thể tiếp tục nhận gradient trong pha 2; việc hợp nhất rồi đóng băng chỉ phù hợp như một cấu hình thí nghiệm loại bỏ thành phần. Bảng \ref{tab:ch2_peft} đặt lựa chọn này trong tương quan với các chiến lược thích nghi khác.

\begin{table}[htbp]
\centering
\small
\renewcommand{\arraystretch}{1.2}
\begin{tabularx}{\textwidth}{L{2.5cm} L{3.0cm} X X}
\toprule
\textbf{Chiến lược} & \textbf{Tham số được cập nhật} & \textbf{Đặc điểm} & \textbf{Mức liên quan với luận văn} \\
\midrule
Tinh chỉnh toàn phần & Toàn bộ backbone và đầu ra & Khả năng thích nghi cao nhưng chi phí huấn luyện lớn & Đối chứng về hiệu năng và tài nguyên. \\
\addlinespace
Dò tuyến tính & Chỉ đầu phân loại tuyến tính & Backbone được giữ cố định; chi phí huấn luyện thấp & Đối chứng cho mức độ hữu ích của biểu diễn chưa thích nghi. \\
\addlinespace
Prompt tuning & Prompt hoặc token có thể học & Giữ phần lớn backbone cố định; hiệu quả phụ thuộc thiết kế prompt & Hướng liên quan nhưng không phải thành phần pipeline chính \cite{PLIP,BiomedCoOp,FPT}. \\
\addlinespace
Adapter / khối nối & Mô-đun nhỏ được chèn thêm & Có tính mô-đun; có thể tăng độ sâu đường suy luận & Tham khảo cho kết nối đa phương thức \cite{FTB}. \\
\addlinespace
LoRA & Hai ma trận hạng thấp cho các lớp được chọn & Giảm số tham số huấn luyện và vẫn thích nghi biểu diễn & Chiến lược được dùng ở cả hai pha \cite{Hu2022LoRA,LCBB}. \\
\bottomrule
\end{tabularx}
\caption{So sánh các chiến lược thích nghi mô hình nền tảng.}
\label{tab:ch2_peft}
\end{table}

\section{Dung hợp đa phương thức và huấn luyện hai pha}

\subsection{Từ phép nối đặc trưng đến chú ý chéo}

Phép nối đặc trưng toàn cục là một baseline đơn giản: hai vector ảnh và văn bản được ghép trước khi đưa vào MLP. Cơ chế này cho phép bộ phân loại sử dụng cả hai nguồn dữ liệu nhưng không mô hình hóa tường minh quan hệ giữa các token. Chú ý chéo cho phép một phương thức truy vấn phương thức còn lại:
\begin{equation}
    \operatorname{Attn}(Q,K,V)
    = \operatorname{softmax}\!\left(\frac{QK^{\top}}{\sqrt{d_k}}\right)V.
\end{equation}
Khi $Q$ được tạo từ ảnh và $(K,V)$ từ văn bản, biểu diễn ảnh được điều kiện hóa bởi token lâm sàng. Chiều ngược lại cho phép biểu diễn văn bản nhận thông tin từ các token ảnh. TieNet \cite{TieNet} và IRENE \cite{IRENE} cho thấy attention có thể được sử dụng để mô hình hóa tương tác đa phương thức; tuy nhiên, chúng không trực tiếp xác nhận rằng chú ý hai chiều luôn tốt hơn phép nối trong mọi bộ dữ liệu.

Pipeline đề xuất vì vậy sử dụng chú ý chéo hai chiều như một giả thuyết thiết kế cần được kiểm chứng. Hai nhánh đầu ra được cộng residual với truy vấn toàn cục, chuẩn hóa và đưa qua MLP dung hợp. Các đối chứng gồm phép nối và chú ý một chiều được sử dụng để tách ảnh hưởng của cơ chế dung hợp khỏi ảnh hưởng của backbone hoặc dữ liệu đầu vào.

\subsection{Vai trò của hai pha huấn luyện}

Pha 1 thích nghi hai bộ mã hóa và mô-đun token không gian bằng Semantic Matching trên cặp ảnh X-quang và bệnh sử; mô-đun dung hợp và đầu phân loại chưa tham gia vào pha này. Mục tiêu của pha 1 là điều chỉnh không gian ảnh--văn bản trước khi tối ưu trực tiếp cho phân loại. Pha 2 tiếp tục cập nhật các LoRA adapter và mô-đun token không gian, đồng thời bắt đầu tối ưu mô-đun chú ý chéo. Cách tách hai pha là một tổng hợp thiết kế từ học tương phản y khoa, PEFT và dung hợp token; nó không được xem là bản sao nguyên vẹn của một công trình đơn lẻ.

Việc tiếp tục tối ưu LoRA trong pha 2 cho phép gradient phân loại tác động đến biểu diễn của hai phương thức trong phạm vi hạng thấp. Ngược lại, nếu hợp nhất và đóng băng LoRA ngay sau pha 1, mô-đun dung hợp phải thích nghi với một backbone cố định. Hai chiến lược này cần được so sánh trên cùng thiết lập để đánh giá vai trò của việc thích nghi liên tục.

\section{Phân loại dựa trên tâm lớp và dữ liệu mất cân bằng}

\subsection{Từ prototypical networks đến tâm lớp thực nghiệm}

Prototypical Networks \cite{Snell2017ProtoNet} biểu diễn mỗi lớp bằng trung bình embedding của support set và phân loại query theo khoảng cách đến các prototype. Phương pháp này thường đi kèm episodic training nhằm mô phỏng tác vụ few-shot. SimpleShot \cite{Wang2019SimpleShot} sử dụng quy tắc láng giềng gần nhất trong thiết lập one-shot và tâm trung bình của từng lớp trong thiết lập multi-shot; nghiên cứu này cũng xem xét phép trừ trung bình và chuẩn hóa $L_2$ mà không cần meta-learning. Vì vậy, biến thể multi-shot của SimpleShot là tiền đề gần hơn cho đầu empirical centroid so với một bộ phân loại kNN trên từng mẫu train.

Các ý tưởng trên tạo cơ sở cho phân loại theo tâm lớp, nhưng cần phân biệt với empirical centroid head của luận văn. Với lớp $k$, tâm lớp được tính từ toàn bộ mẫu huấn luyện hiện có:
\begin{equation}
    c_k = \frac{1}{N_k}\sum_{i:y_i=k} z_i.
\end{equation}
Sau khi chuẩn hóa $L_2$, logit của mẫu $z$ đối với lớp $k$ được tính bằng độ tương đồng cosine có hệ số tỉ lệ $\alpha$:
\begin{equation}
    \ell_k(z) = \alpha\,
    \frac{z^{\top}c_k}{\lVert z\rVert_2\lVert c_k\rVert_2}.
\end{equation}
Các centroid là thống kê của tập huấn luyện, được lưu như buffer và không được optimizer cập nhật như tham số tự do. Khi encoder thay đổi, centroid được ước lượng lại từ embedding train. Vì vậy, pipeline không sử dụng Prototypical Loss dạng pull--push và cũng không cần episodic training để học prototype. Bảng \ref{tab:ch2_centroid} làm rõ sự khác biệt giữa centroid thực nghiệm và các dạng đại diện lớp có nhận gradient.

\begin{table}[htbp]
\centering
\small
\renewcommand{\arraystretch}{1.2}
\begin{tabularx}{\textwidth}{L{2.7cm} L{3.1cm} X X}
\toprule
\textbf{Đầu phân loại} & \textbf{Cách xác định đại diện lớp} & \textbf{Cách huấn luyện} & \textbf{Nhận xét} \\
\midrule
Linear head & Trọng số học được cho từng lớp & Tối ưu bằng gradient cùng encoder hoặc riêng đầu ra & Linh hoạt nhưng thêm tham số phân loại. \\
\addlinespace
Prototype học được & Vector lớp là tham số tự do & Tối ưu trực tiếp bằng loss phân loại hoặc metric loss & Có thể lệch khỏi trung bình embedding train nếu không có ràng buộc. \\
\addlinespace
Prototypical Networks \cite{Snell2017ProtoNet} & Trung bình support set trong từng episode & Episodic training trên support--query & Nền tảng của phân loại theo prototype trong few-shot. \\
\addlinespace
SimpleShot \cite{Wang2019SimpleShot} & Mẫu support trong one-shot hoặc trung bình lớp trong multi-shot & Không dùng meta-learning; biến đổi đặc trưng trước phân loại láng giềng & Tiền đề gần với phân loại theo tâm lớp trên đặc trưng chuẩn hóa. \\
\addlinespace
Empirical centroid & Trung bình toàn bộ embedding train của mỗi lớp & Centroid cố định trong epoch và được ước lượng lại khi encoder đổi & Phù hợp với mục tiêu đầu phân loại phi tham số của luận văn. \\
\bottomrule
\end{tabularx}
\caption{Phân biệt các dạng đầu phân loại dựa trên trọng số và tâm lớp.}
\label{tab:ch2_centroid}
\end{table}

\subsection{Cross-Entropy có trọng số theo số mẫu hiệu dụng}

Khi số mẫu giữa các lớp chênh lệch, Cross-Entropy không trọng số có thể làm gradient bị chi phối nhiều hơn bởi các lớp phổ biến. Cui và cộng sự \cite{Cui2019ClassBalanced} đề xuất số mẫu hiệu dụng của lớp có $N_k$ quan sát:
\begin{equation}
    E_k = \frac{1-\beta^{N_k}}{1-\beta}, \qquad 0 \leq \beta < 1,
\end{equation}
và sử dụng trọng số tỉ lệ nghịch với $E_k$. Loss phân loại trong pha 2 có dạng:
\begin{equation}
    \mathcal{L}_{\mathrm{cls}}
    = -\frac{1}{N}\sum_{i=1}^{N}w_{y_i}
      \log p(y_i\mid z_i),
\end{equation}
trong đó các trọng số được chuẩn hóa trên những lớp hiện diện. Phương pháp này điều chỉnh đóng góp của lớp trong hàm mục tiêu mà không tạo thêm mẫu tổng hợp. Hiệu quả thực tế vẫn phụ thuộc vào giá trị $\beta$, số mẫu và chất lượng biểu diễn, do đó cần được báo cáo cùng các chỉ số macro và kết quả theo lớp.

\section{Phát hiện dữ liệu ngoài phân phối}

\subsection{Bài toán và nguyên tắc đánh giá}

Phát hiện OOD nhằm gán một điểm bất thường cho mẫu không phù hợp với phân phối huấn luyện. Dữ liệu thuộc phân phối (In-Distribution -- ID) nhưng khó vẫn có thể bị phân loại sai, trong khi một mẫu OOD đôi khi vẫn nhận logit cao ở một lớp đã biết. Vì vậy, xác suất dự đoán không nên được xem như bảo đảm về độ tin cậy.

Với điểm OOD $s(x)$ được quy ước càng lớn càng bất thường, ngưỡng từ chối phải được lựa chọn trên tập hiệu chỉnh độc lập với tập kiểm thử. Việc dùng mẫu OOD test để chọn phương pháp, tham số hoặc ngưỡng có thể làm sai lệch đánh giá. Pipeline khảo sát song song nhóm điểm dựa trên đầu ra và nhóm điểm dựa trên không gian đặc trưng; kết quả cuối cùng được báo cáo bằng AUROC, AUPR-Out và FPR@95TPR theo các kịch bản OOD được xác định trước. Trong quy ước đánh giá của luận văn, FPR@95TPR là tỷ lệ mẫu OOD bị chấp nhận nhầm như ID khi tỷ lệ chấp nhận đúng mẫu ID được giữ ở mức 95\%.

\subsection{Điểm dựa trên đầu ra mô hình}

Maximum Softmax Probability (MSP) \cite{Hendrycks2017MSP} sử dụng xác suất Softmax lớn nhất làm baseline; trong quy ước điểm bất thường có thể dùng $1-\max_k p_k(x)$. Entropy đánh giá độ phân tán của toàn bộ phân phối dự đoán. Maximum Logit \cite{Hendrycks2022MaxLogit} sử dụng logit lớn nhất trước Softmax và có thể giảm bớt ảnh hưởng của phép chuẩn hóa giữa các lớp.

Energy score \cite{Liu2020EnergyOOD} được tính từ vector logit $f(x)$:
\begin{equation}
    E(x)=-T\log\sum_{k=1}^{K}
    \exp\!\left(\frac{f_k(x)}{T}\right),
\end{equation}
trong đó $T$ là nhiệt độ. Khi sử dụng quy ước trên, giá trị năng lượng lớn hơn được xem là bằng chứng OOD mạnh hơn. Các điểm đầu ra có ưu điểm là không cần lưu toàn bộ embedding train, nhưng có thể chịu ảnh hưởng của mức hiệu chuẩn và hình học của đầu phân loại.

\subsection{Điểm dựa trên không gian đặc trưng}

Lee và cộng sự \cite{Lee2018MahalanobisOOD} mô hình hóa đặc trưng theo phân bố Gaussian có điều kiện theo lớp và dùng khoảng cách Mahalanobis làm điểm tin cậy. Với embedding đã chuẩn hóa, tâm lớp $c_k$ và ma trận hiệp phương sai dùng chung $\Sigma$ ước lượng từ phần dư trong lớp, biến thể được sử dụng trong luận văn có dạng:
\begin{equation}
    s_{\mathrm{Maha}}(z)=
    \min_k (z-c_k)^{\top}\Sigma^{-1}(z-c_k).
\end{equation}
Cách viết này khác với giả định một ma trận $\Sigma_k$ riêng cho từng lớp. Hiệp phương sai dùng chung giúp giảm số lượng đại lượng cần ước lượng khi một số lớp có ít mẫu; pipeline còn sử dụng shrinkage Ledoit--Wolf \cite{LedoitWolf2004} khi ước lượng để cải thiện tính ổn định số.

Deep k-Nearest Neighbors \cite{Sun2022KNN} sử dụng khoảng cách đến các embedding train gần nhất và không cần giả định Gaussian. Trong luận văn, các vector được chuẩn hóa và khoảng cách cosine trung bình tới $k$ láng giềng gần nhất được dùng làm điểm OOD. Phương pháp này linh hoạt hơn về hình dạng phân bố nhưng cần lưu tập đặc trưng tham chiếu và có chi phí truy vấn theo kích thước tập train nếu không dùng chỉ mục gần đúng.

Đối với mô hình ngôn ngữ -- thị giác, MCM \cite{Ming2022MCM} sử dụng độ tương đồng giữa ảnh và các khái niệm văn bản để phát hiện OOD zero-shot. Theo hướng tương tự, mã nguồn hỗ trợ khoảng cách đến các mốc văn bản của lớp ID. Tuy nhiên, phương pháp này chưa thuộc protocol OOD được khóa cho thí nghiệm chính, nên được xem là hướng mở rộng; muốn báo cáo như một kết quả kiểm định, cần bổ sung nó vào protocol và xác định trước cách xây dựng prompt, tập hiệu chỉnh và siêu tham số. Bảng \ref{tab:ch2_ood} tổng hợp tín hiệu đầu vào và giả định của từng nhóm điểm OOD.

\begin{table}[htbp]
\centering
\small
\renewcommand{\arraystretch}{1.2}
\begin{tabularx}{\textwidth}{L{2.3cm} L{2.4cm} X X}
\toprule
\textbf{Phương pháp} & \textbf{Tín hiệu sử dụng} & \textbf{Giả định / đặc điểm} & \textbf{Vai trò trong đánh giá} \\
\midrule
MSP \cite{Hendrycks2017MSP} & Xác suất Softmax & Baseline đơn giản; nhạy với dự đoán quá tự tin & So sánh đầu ra hậu xử lý. \\
\addlinespace
Entropy & Phân phối xác suất & Xem xét độ phân tán trên tất cả lớp & Baseline bổ sung, không cần mô hình mới. \\
\addlinespace
Maximum Logit \cite{Hendrycks2022MaxLogit} & Logit lớn nhất & Không chuẩn hóa qua tổng Softmax & So sánh với MSP và Energy. \\
\addlinespace
Energy \cite{Liu2020EnergyOOD} & Toàn bộ vector logit & Điểm vô hướng hậu xử lý có nhiệt độ & Phương pháp đầu ra chính được khảo sát. \\
\addlinespace
Mahalanobis \cite{Lee2018MahalanobisOOD} & Tâm lớp và hiệp phương sai & Giả định hình học Gaussian trong feature space & Điểm mật độ tham số trên embedding. \\
\addlinespace
Cosine-kNN \cite{Sun2022KNN} & Láng giềng trong tập train & Phi tham số, không giả định dạng cụm & Điểm khoảng cách bổ sung cho Mahalanobis. \\
\addlinespace
Mốc văn bản \cite{Ming2022MCM} & Tương đồng ảnh--khái niệm & Phụ thuộc prompt và khả năng căn chỉnh VLM & Hướng mở rộng đã được mã nguồn hỗ trợ nhưng chưa nằm trong protocol chính. \\
\bottomrule
\end{tabularx}
\caption{Các nhóm điểm OOD có liên quan đến pipeline đánh giá.}
\label{tab:ch2_ood}
\end{table}

\section{Trực quan hóa và phân tích mức đóng góp}

Grad-CAM \cite{Selvaraju2017GradCAM} tạo bản đồ định vị theo lớp từ gradient và activation của mạng tích chập, qua đó cung cấp một baseline phổ biến cho trực quan hóa vùng ảnh. Với Transformer, attention weight thường được sử dụng như một tín hiệu để khảo sát quan hệ token. Tuy nhiên, attention ở nhiều tầng làm thông tin nguồn bị trộn; attention rollout và attention flow \cite{Abnar2020AttentionFlow} được đề xuất để tổng hợp đường truyền attention thay vì chỉ đọc một ma trận đơn lẻ.

Việc diễn giải attention cần được thực hiện thận trọng. Jain và Wallace \cite{Jain2019AttentionNotExplanation} chỉ ra rằng attention weight không mặc nhiên tương ứng với mức quan trọng nhân quả của đầu vào. Vì vậy, bản đồ attention trong luận văn được xem là công cụ phân tích hành vi và hỗ trợ kiểm tra vị trí mô hình tập trung, không phải bằng chứng duy nhất giải thích quyết định hoặc thay thế đánh giá của bác sĩ.

Trong XBone-Net, attention từ nhánh văn bản đến token ảnh có thể được tổng hợp và chiếu ngược về tọa độ các vùng nguồn. Nếu một tập đánh giá có chú thích tổn thương, mức tập trung có thể được định lượng bằng tỷ lệ khối lượng attention nằm trong vùng tổn thương, tỷ lệ diện tích vùng tổn thương và tỷ số giữa hai đại lượng. So sánh với tỷ lệ diện tích giúp hạn chế kết luận dựa trên bản đồ nhiệt trực quan đơn thuần. Protocol CTCH hiện không báo cáo định vị không gian khi chưa có chú thích vùng tổn thương; do đó, công cụ này chỉ thể hiện khả năng hỗ trợ phân tích có điều kiện, không đồng nghĩa đã có kết quả định lượng định vị trong thí nghiệm chính. Các phép kiểm tra deletion/insertion hoặc gradient-based attribution có thể được dùng như phân tích bổ sung để đánh giá mức liên hệ giữa bản đồ và đầu ra mô hình. Bảng \ref{tab:ch2_explainability} phân biệt giá trị sử dụng và giới hạn của các tín hiệu trực quan hóa.

\begin{table}[htbp]
\centering
\small
\renewcommand{\arraystretch}{1.2}
\begin{tabularx}{\textwidth}{L{2.5cm} L{2.7cm} X X}
\toprule
\textbf{Hướng phân tích} & \textbf{Tín hiệu} & \textbf{Giá trị sử dụng} & \textbf{Giới hạn} \\
\midrule
Grad-CAM \cite{Selvaraju2017GradCAM} & Gradient và activation & Bản đồ theo lớp cho mạng tích chập & Không trùng trực tiếp với attention của mô hình đa phương thức. \\
\addlinespace
Attention trực tiếp & Trọng số giữa query và token & Thuận tiện để khảo sát tương tác ảnh--văn bản & Không mặc nhiên phản ánh đóng góp nhân quả \cite{Jain2019AttentionNotExplanation}. \\
\addlinespace
Attention rollout / flow \cite{Abnar2020AttentionFlow} & Attention qua nhiều tầng & Theo dõi sự truyền thông tin giữa các lớp Transformer & Vẫn là phép xấp xỉ hậu xử lý. \\
\addlinespace
Định lượng trên vùng tổn thương & Attention đã chiếu về ảnh và chú thích vùng & Cho phép so sánh với tỷ lệ diện tích khi có annotation & Phụ thuộc chất lượng chú thích và phép chiếu token--pixel; chưa báo cáo cho CTCH khi thiếu chú thích vùng. \\
\bottomrule
\end{tabularx}
\caption{Các hướng trực quan hóa và giới hạn diễn giải.}
\label{tab:ch2_explainability}
\end{table}

\section{Bàn luận và khoảng trống nghiên cứu}

\subsection{Những điểm đã được hỗ trợ bởi nghiên cứu trước}

Tổng quan cho thấy từng thành phần của pipeline đều có tiền đề trong các nhánh nghiên cứu hiện có. MedCLIP và BiomedCLIP hỗ trợ hướng căn chỉnh rồi thích nghi mô hình ngôn ngữ -- thị giác y khoa. TieNet, IRENE và FTB cung cấp các ví dụ về tương tác ảnh--văn bản. LoRA tạo cơ sở cho thích nghi hạng thấp. Prototypical Networks và SimpleShot cung cấp góc nhìn về phân loại theo đại diện lớp, trong khi class-balanced loss điều chỉnh mức đóng góp không đồng đều giữa các lớp. Các phương pháp Mahalanobis, kNN, Energy và MCM cung cấp những điểm OOD có giả định bổ sung cho nhau.

Tuy vậy, kết quả từ các công trình này không trực tiếp chứng minh hiệu quả của tổ hợp cụ thể trên dữ liệu u xương. Phần lớn nghiên cứu sử dụng X-quang ngực, CT, ảnh mô bệnh học hoặc bộ dữ liệu tự nhiên; loại văn bản và phân bố lớp cũng khác nhau. Do đó, mỗi quyết định của XBone-Net cần được đánh giá bằng baseline và ablation trên cùng cách chia dữ liệu.

\subsection{Khoảng trống đối với bài toán X-quang xương}

Thứ nhất, nhiều công trình được khảo sát ở trên sử dụng báo cáo X-quang, trong khi điều kiện dự đoán của đề tài yêu cầu sử dụng thông tin có trước kết luận đọc ảnh. Điều này đặt ra nhu cầu kiểm soát nguồn văn bản và rò rỉ nhãn. Thứ hai, ảnh X-quang xương có thể chứa tổn thương cục bộ nhưng backbone chỉ nhận kích thước cố định, tạo sự đánh đổi giữa ngữ cảnh toàn cục và chi tiết. Thứ ba, dữ liệu hạn chế và mất cân bằng làm tăng nhu cầu về PEFT, đầu phân loại ít tham số và hàm mất mát có trọng số. Thứ tư, một hệ thống phân loại đóng cần cơ chế từ chối các mẫu khác phân phối. Cuối cùng, bản đồ attention cần được định lượng và trình bày với giới hạn diễn giải rõ ràng. Bảng \ref{tab:ch2_gap_mapping} ánh xạ các khoảng trống này tới thành phần thiết kế và nội dung cần kiểm chứng.

\begin{table}[htbp]
\centering
\small
\renewcommand{\arraystretch}{1.2}
\begin{tabularx}{\textwidth}{L{3.0cm} L{3.6cm} X X}
\toprule
\textbf{Vấn đề} & \textbf{Tiền đề nghiên cứu} & \textbf{Lựa chọn trong XBone-Net} & \textbf{Nội dung cần kiểm chứng} \\
\midrule
Kết hợp bối cảnh lâm sàng & Mei, TieNet và IRENE \cite{XMei,TieNet,IRENE} & Ảnh X-quang và bệnh sử trước khi có kết luận đọc ảnh ở cả hai pha & So sánh đa phương thức với ảnh đơn thuần; đánh giá ảnh hưởng của báo cáo đọc ảnh trong pha căn chỉnh. \\
\addlinespace
Căn chỉnh miền y sinh & MedCLIP và BiomedCLIP \cite{MedCLIP,BiomedCLIP} & Semantic Matching ở pha 1 trên backbone BiomedCLIP & So sánh với baseline không có pha căn chỉnh hoặc mục tiêu khác. \\
\addlinespace
Ảnh độ phân giải cao & FPT, Perceiver và Flamingo \cite{FPT,Jaegle2021Perceiver,Alayrac2022Flamingo} & Global view, local view và token có tọa độ với ngân sách cố định & Ablation nhánh độ phân giải cao; cần cấu hình bổ sung nếu muốn cô lập riêng tác động của tọa độ. \\
\addlinespace
Thích nghi ít tham số & LoRA và LCBB \cite{Hu2022LoRA,LCBB} & LoRA tiếp tục được tối ưu qua hai pha & So sánh với tinh chỉnh toàn phần và chiến lược đóng băng sau pha 1. \\
\addlinespace
Phân loại mất cân bằng & ProtoNet, SimpleShot và effective number \cite{Snell2017ProtoNet,Wang2019SimpleShot,Cui2019ClassBalanced} & Empirical centroid và weighted Cross-Entropy & Báo cáo Macro F1, kết quả theo lớp và ablation đầu phân loại/loss. \\
\addlinespace
Phát hiện OOD & MSP, Energy, Mahalanobis, kNN và MCM \cite{Hendrycks2017MSP,Liu2020EnergyOOD,Lee2018MahalanobisOOD,Sun2022KNN,Ming2022MCM} & Sáu điểm trong protocol chính; mốc văn bản là hướng mở rộng & So sánh trên các kịch bản xác định trước, hiệu chỉnh ngưỡng bằng validation và không dùng test để lựa chọn. \\
\addlinespace
Phân tích vùng chú ý & Grad-CAM và attention flow \cite{Selvaraju2017GradCAM,Abnar2020AttentionFlow} & Chiếu attention về ảnh; định lượng chỉ khi có chú thích vùng & Phân tích định tính trên CTCH và nêu rõ giới hạn; không báo cáo định vị nếu thiếu annotation. \\
\bottomrule
\end{tabularx}
\caption{Liên hệ giữa khoảng trống nghiên cứu, lựa chọn thiết kế và nội dung đánh giá.}
\label{tab:ch2_gap_mapping}
\end{table}

\subsection{Cơ sở hình thành pipeline đề xuất}

Từ các phân tích trên, pipeline được xây dựng theo hai pha. Pha 1 thích nghi BiomedCLIP bằng LoRA và Semantic Matching trên cặp ảnh X-quang--bệnh sử lâm sàng; mô-đun token không gian thuộc nhánh thị giác cũng nhận gradient trong pha này. Pha 2 tiếp tục cập nhật LoRA và mô-đun token không gian, đồng thời tối ưu chú ý chéo hai chiều; empirical centroid được ước lượng từ embedding train và Cross-Entropy được tái trọng số theo số mẫu hiệu dụng. Sau huấn luyện, Mahalanobis và kNN được khớp trên đặc trưng train, còn các điểm dựa trên logit được tính trực tiếp từ đầu ra mô hình; ngưỡng của từng phương pháp được hiệu chỉnh trên validation. Attention được trích xuất như một tín hiệu phân tích, chiếu về ảnh nguồn và đánh giá trên vùng tổn thương khi có chú thích.

Thiết kế này không giả định rằng mọi thành phần đều mang lại cải thiện trong mọi điều kiện. Vai trò của từng thành phần được chuyển thành các câu hỏi nghiên cứu và được kiểm tra bằng mốc đối chứng, thí nghiệm loại bỏ thành phần, chỉ số macro, đánh giá OOD và phân tích attention; chỉ số định vị chỉ được tính khi có chú thích vùng phù hợp. Kiến trúc, hàm mục tiêu và quy trình đánh giá cụ thể được trình bày trong Chương 3.
